% ---------------------------------------------------------------------------
% The capability-state boundary theorem for the crumbling-wall quorum system.
% Drop-in LaTeX.  Requires amsthm/amsmath and the usual theorem environments:
%   \newtheorem{theorem}{Theorem}[section]
%   \newtheorem{lemma}[theorem]{Lemma}
%   \newtheorem{proposition}[theorem]{Proposition}
%   \newtheorem{corollary}[theorem]{Corollary}
%   \theoremstyle{definition} \newtheorem{definition}[theorem]{Definition}
%   \theoremstyle{remark}     \newtheorem{remark}[theorem]{Remark}
% ---------------------------------------------------------------------------

\subsection{Capability states and the preemption hazard}

\begin{definition}[Wall]
\label{def:wall}
A \emph{wall} is a finite node set $V$ together with an ordered list of
\emph{tiers} $T_0,\dots,T_{m-1}$ with $T_j\subseteq V$ and
$\bigcup_j T_j=V$.  The last tier $E:=T_{m-1}$ is the \emph{anchor}
(``Earth''); write $n_E:=|E|$.  Tiers are listed slow-to-fast, matching the
code convention of \texttt{CrumblingWallQuorum}; the paper's indexing is the
reverse.  Fix a \emph{Phase~2 threshold} $k$ and set
\[
  a \;:=\; n_E-k+1 ,
\]
the minimum number of anchor nodes any Phase~1 quorum must contain
(\texttt{min\_earth\_in\_q1}).  For a respondent set $R\subseteq V$ and an
initiator tier $i\in\{0,\dots,m-1\}$ define
\begin{align}
  \Phi_1^{\,i}(R) &\;:\Longleftrightarrow\;
    \Bigl(\forall j,\; i\le j\le m-1:\; R\cap T_j\neq\emptyset\Bigr)
    \;\wedge\; |R\cap E|\ \ge\ a ,
    \label{eq:p1}\\[2pt]
  \Phi_2(R) &\;:\Longleftrightarrow\; |R\cap E|\ \ge\ k .
    \label{eq:p2}
\end{align}
The \emph{capability state} of a tier-$i$ proposer whose reachable set is $R$
is the pair $\chi^i(R):=\bigl(\Phi_1^{\,i}(R),\Phi_2(R)\bigr)\in\{0,1\}^2$.
The state $(1,0)$ --- able to complete Phase~1, hence able to raise the ballot
high-water mark and preempt incumbents, but unable to complete Phase~2, hence
unable to commit --- is the \emph{preemption hazard}.
\end{definition}

Equations \eqref{eq:p1}--\eqref{eq:p2} are a literal transcription of
\texttt{CrumblingWallQuorum.is\_phase1\_quorum} and
\texttt{is\_phase2\_quorum} in \texttt{quorums.py}.

\begin{lemma}[Monotonicity]
\label{lem:mono}
For every $i$, both $\Phi_1^{\,i}$ and $\Phi_2$ are monotone: if
$R\subseteq R'$ and $\Phi(R)$ then $\Phi(R')$.  Consequently a tier-$i$
proposer can \emph{assemble} a Phase~$p$ quorum from a reachable set $R$ if
and only if $R$ itself satisfies $\Phi_p$.
\end{lemma}

\begin{proof}
Each conjunct of \eqref{eq:p1} and \eqref{eq:p2} is of the form
``$R\cap S\neq\emptyset$'' or ``$|R\cap S|\ge c$'' for a fixed $S$ and
constant $c$, and both are preserved under enlarging $R$.  For the
consequence: if some $Q\subseteq R$ satisfies $\Phi_p$ then so does $R$ by
monotonicity; conversely $R\subseteq R$.
\end{proof}

Lemma~\ref{lem:mono} is what licenses the phrase ``reachable capability
state'': we may quantify over respondent sets rather than over quorums.

\subsection{The boundary}

\begin{theorem}[Capability boundary]
\label{thm:boundary}
Let a wall be given with anchor $E$, $n_E\ge 1$, and threshold $k$ with
\begin{description}
  \item[(W1)] $1\le k\le n_E$ \quad(equivalently $1\le a\le n_E$), and
  \item[(W2)] $T_j\setminus E\neq\emptyset$ for every $j$ with
    $0\le j\le m-2$ \quad(every non-anchor tier owns at least one node the
    anchor does not).
\end{description}
Then for every initiator tier $i$,
\[
  \bigl(\exists R\subseteq V:\ \chi^i(R)=(1,0)\bigr)
  \quad\Longleftrightarrow\quad
  a<k
  \quad\Longleftrightarrow\quad
  n_E-k+1<k
  \quad\Longleftrightarrow\quad
  k>\Bigl\lceil \tfrac{n_E}{2}\Bigr\rceil .
\]
Equivalently: the hazard is unreachable if and only if
$n_E-k+1\ge k$, i.e.\ $k\le\lceil n_E/2\rceil$.
\end{theorem}

\begin{proof}
The arithmetic equivalences first.  Since $a=n_E-k+1$ we have
$a\ge k\iff n_E+1\ge 2k\iff k\le\lfloor (n_E+1)/2\rfloor$, and
$\lfloor (n_E+1)/2\rfloor=\lceil n_E/2\rceil$ for every integer $n_E$.
Note also the identity $a+k=n_E+1$: the two thresholds partition
$n_E+1$, and $a\ge k$ says exactly that Phase~2 receives the smaller half.

\smallskip
\emph{Sufficiency ($a\ge k$ $\Rightarrow$ no hazard).}
Suppose $\Phi_1^{\,i}(R)$ holds.  Its second conjunct gives
$|R\cap E|\ge a\ge k$, which is precisely $\Phi_2(R)$.  Hence
$\chi^i(R)\in\{(0,0),(0,1),(1,1)\}$ and $(1,0)$ is unreachable.  This
direction uses \emph{only} the anchor-count conjunct: it holds for every $i$,
every $m$, every choice of tier sizes, and with or without (W1) and (W2).

\smallskip
\emph{Necessity ($a<k$ $\Rightarrow$ hazard).}
Assume $a\le k-1$.  By (W1), $1\le a\le n_E$, so we may choose $A\subseteq E$
with $|A|=a$.  By (W2), for each $j$ with $i\le j\le m-2$ choose a witness
$w_j\in T_j\setminus E$.  Put
\[
  R \;:=\; A\ \cup\ \{\,w_j \;:\; i\le j\le m-2\,\}, \qquad
  |R| \;=\; (m-1-i)+a .
\]
\emph{$R$ satisfies $\Phi_1^{\,i}$.}  For $i\le j\le m-2$ we have
$w_j\in R\cap T_j$, so that tier is hit.  For $j=m-1$ we have
$A\subseteq R\cap E$ and $|A|=a\ge1$, so the anchor is hit.  Finally
$|R\cap E|\ge|A|=a$.
\emph{$R$ violates $\Phi_2$.}  Every $w_j$ lies outside $E$, so
$R\cap E=A$ exactly and $|R\cap E|=a\le k-1<k$.
Hence $\chi^i(R)=(1,0)$.
\end{proof}

\begin{remark}[What the theorem does \emph{not} depend on]
\label{rem:independence}
Three quantities that one might expect to matter do not appear in the
criterion.
\begin{enumerate}
\item \textbf{The number of tiers $m$.}  It appears only in the size
  $(m-1-i)+a$ of the witness set, never in the predicate that decides
  reachability.  The per-tier conjuncts of \eqref{eq:p1} are
  \emph{hitting} constraints on sets disjoint from $E$ (by (W2)); satisfying
  them costs nodes, but costs zero \emph{anchor} nodes, so they are inert with
  respect to the only quantity $\Phi_2$ measures.
\item \textbf{The sizes of the non-anchor tiers.}  Only emptiness matters
  ((W2)); a tier of size $1$ and a tier of size $10^6$ impose the same
  constraint ``contribute one node, none of them in $E$''.  Tier sizes do
  change the \emph{number} of hazard states --- each non-anchor tier
  multiplies the count of hazardous respondent sets --- but not their
  existence.
\item \textbf{The initiator tier $i$.}  Sufficiency is uniform in $i$ because
  the anchor-count conjunct of \eqref{eq:p1} is the same for all $i$.
  Necessity for tier $i$ needs (W2) only for $j\ge i$, which is implied by
  (W2) for $j\ge 0$.  So all tiers cross the boundary at the same $k$: the
  wall has one boundary, not one per tier.
\end{enumerate}
The criterion is therefore a statement about a single coordinate --- the
number of anchor nodes in $R$ --- projected out of an otherwise
higher-dimensional constraint system.
\end{remark}

\begin{remark}[Self-inclusion]
Whether a proposer counts itself among the respondents is immaterial:
adding a fixed node to $R$ preserves $\Phi_1^{\,i}$ by
Lemma~\ref{lem:mono}, and in the necessity construction one may take the
proposer's own node as the witness $w_i$ (if $i\le m-2$) or as a member of
$A$ (if $i=m-1$) without changing $|R\cap E|$ or the argument.
\end{remark}

\subsection{The flat-threshold companion}

\begin{corollary}[Threshold quorums]
\label{cor:flat}
Let $\Phi_1(R)\iff|R|\ge q_1$ and $\Phi_2(R)\iff|R|\ge q_2$ over $n$ nodes,
with $1\le q_1,q_2\le n$.  Then $(1,0)$ is unreachable if and only if
$q_1\ge q_2$.
\end{corollary}

\begin{proof}
If $q_1\ge q_2$ then $|R|\ge q_1\ge q_2$, so $\Phi_1(R)\Rightarrow\Phi_2(R)$.
If $q_1<q_2$, any $R$ with $|R|=q_1\ (\le n)$ satisfies $\Phi_1$ and violates
$\Phi_2$.
\end{proof}

Corollary~\ref{cor:flat} is Theorem~\ref{thm:boundary} with $m=1$ (no
non-anchor tiers, so (W2) is vacuous), $E=V$, $a=q_1$, $k=q_2$.  The two
results are one lemma --- \emph{two thresholds on a common set are
comparable, and the hazard is exactly the gap between them} --- and the
crumbling wall is the flat rule applied \emph{inside the anchor}: the
Flexible Paxos condition $q_1+q_2>n$ instantiated on $E$ reads
$a+k>n_E$, and the construction takes the minimum such $a=n_E-k+1$.  Under
that choice $a\ge k$ is literally $q_1\ge q_2$ restricted to the anchor.
The hazard is reachable exactly when Phase~1 is \emph{cheaper} than Phase~2
on the anchor.

\subsection{Degenerate walls: a sharp form}

Hypotheses (W1)--(W2) are not cosmetic; Section~\ref{sec:cx} exhibits walls
violating them where the ``only if'' direction of
Theorem~\ref{thm:boundary} fails.  The following proposition is exact for all
walls, degenerate or not, and shows precisely how the boundary deforms.

\begin{definition}
For an initiator tier $i$ let
$\mathcal F_i:=\{\,T_j : i\le j\le m-2,\ T_j\subseteq E\,\}$ be the family of
non-anchor tiers swallowed by the anchor, and let $h_i$ be the minimum size of
a set $S\subseteq E$ meeting every member of $\mathcal F_i$ (the
\emph{hitting number}), with $h_i:=0$ if $\mathcal F_i=\emptyset$ and
$h_i:=\infty$ if some member of $\mathcal F_i$ is empty.  Put
\[
  c_i \;:=\;
  \begin{cases}
    \infty, & \text{if } T_j=\emptyset \text{ for some } i\le j\le m-1,\\
    \max\{a,\ h_i,\ 1\}, & \text{otherwise.}
  \end{cases}
\]
\end{definition}

\begin{proposition}[Sharp boundary]
\label{prop:sharp}
For every wall, every $k\in\mathbb Z$ and every $i$,
\[
  \bigl(\exists R:\ \chi^i(R)=(1,0)\bigr)
  \quad\Longleftrightarrow\quad
  c_i<k .
\]
\end{proposition}

\begin{proof}
($\Rightarrow$) Let $\Phi_1^{\,i}(R)$ and $\neg\Phi_2(R)$.  Each tier
$T_j$, $i\le j\le m-1$, meets $R$, hence is nonempty, so the first case of
the definition of $c_i$ does not apply.  Write $D:=R\cap E$.  Then
$|D|\ge a$; $D\neq\emptyset$ because the anchor conjunct $R\cap T_{m-1}
\neq\emptyset$ holds, so $|D|\ge1$; and for $T_j\in\mathcal F_i$ we have
$\emptyset\neq R\cap T_j\subseteq R\cap E=D$, so $D$ hits $\mathcal F_i$ and
$|D|\ge h_i$.  Therefore $|D|\ge\max\{a,h_i,1\}=c_i$.  Since
$\neg\Phi_2(R)$ gives $|D|\le k-1$, we get $c_i\le k-1<k$.

($\Leftarrow$) Let $c_i<k$; in particular $c_i$ is finite, so every $T_j$
with $i\le j\le m-1$ is nonempty and $E\neq\emptyset$.  Choose a minimum
hitting set $S\subseteq E$ for $\mathcal F_i$, $|S|=h_i\le c_i$.  We claim
$c_i\le n_E$: if $k\le n_E$ then $c_i\le k-1<n_E$; if $k>n_E$ then $a\le 0$
and $c_i=\max\{h_i,1\}\le n_E$ since $h_i\le n_E$ and $n_E\ge1$.  Hence we may
extend $S$ to $A\subseteq E$ with $|A|=c_i$.  For each $j$ with
$i\le j\le m-2$ and $T_j\not\subseteq E$ choose $w_j\in T_j\setminus E$, and
set $R:=A\cup\{w_j\}$.  Every tier $T_j$, $i\le j\le m-2$, is met: by $w_j$ if
$T_j\not\subseteq E$, by $S\subseteq A$ otherwise.  The anchor is met because
$|A|=c_i\ge1$.  Also $R\cap E=A$, so $|R\cap E|=c_i\ge a$, giving
$\Phi_1^{\,i}(R)$; and $|R\cap E|=c_i<k$, giving $\neg\Phi_2(R)$.
\end{proof}

\begin{corollary}
Under (W1) and (W2), $\mathcal F_i=\emptyset$ so $h_i=0$, every tier from $i$
down is nonempty, and $a\ge1$; hence $c_i=a$ and
Proposition~\ref{prop:sharp} reduces to Theorem~\ref{thm:boundary}.
\end{corollary}

\subsection{Sharpness: how each hypothesis fails}
\label{sec:cx}

Each item below breaks the ``only if'' direction of
Theorem~\ref{thm:boundary} (the ``if'' direction is unconditional).

\begin{enumerate}
\item \textbf{An empty tier} (violates (W2); \emph{accepted} by the shipped
  constructor).  Tiers $\langle\emptyset,\{u\},\{v\},E\rangle$ with
  $n_E=5$, $k=5$, so $a=1<k$.  For $i=0$ the read-down conjunct demands a
  respondent from the empty tier, so $\Phi_1^{\,0}$ is unsatisfiable and
  \emph{no} capability state other than $(0,\cdot)$ is reachable: the hazard
  is absent for the wrong reason.  Proposition~\ref{prop:sharp} agrees:
  $c_0=\infty$.  For $i\ge1$ the hazard returns.
\item \textbf{A tier contained in the anchor} (violates (W2); \emph{rejected}
  by the shipped constructor's disjointness check).  Tiers
  $\langle\{e_1\},\{e_2\},\{e_3\},\{e_4\},E\rangle$ with
  $E=\{e_1,\dots,e_5\}$ and $k=4$, so $a=2<k$.  A tier-$0$ Phase~1 quorum
  must contain $e_1,\dots,e_4$, i.e.\ four anchor nodes, which already
  satisfies $\Phi_2$; the hazard is unreachable at $i=0$ despite $a<k$.
  Here $h_0=4$ and $c_0=4\not<4$.  For $i\ge1$, $h_i\le3<4$ and the hazard
  reappears --- the one configuration in which the boundary genuinely
  differs from tier to tier.
\item \textbf{Partial overlap is harmless.}  If each non-anchor tier retains
  one node outside $E$ --- e.g.\ $\langle\{e_1,x_1\},\dots\rangle$ --- then
  $\mathcal F_i=\emptyset$, $c_i=a$, and the boundary is restored.  The
  operative hypothesis is $T_j\setminus E\neq\emptyset$, which is
  \emph{weaker} than the pairwise disjointness the implementation enforces:
  overlaps \emph{among non-anchor tiers} are irrelevant, since witnesses may
  coincide.
\item \textbf{$k$ outside $[1,n_E]$} (violates (W1)) does \emph{not} break the
  criterion.  For $k>n_E$ we get $a\le0$, $\Phi_2$ is unsatisfiable, and every
  Phase~1 quorum is hazardous --- consistent with $a<k$.  For $k\le0$,
  $\Phi_2$ is universally true and the hazard is impossible --- consistent
  with $a\ge k$.  (W1) is needed only to make the specific witness set of
  Theorem~\ref{thm:boundary} well defined;
  Proposition~\ref{prop:sharp} covers the rest.
\item \textbf{$E=\emptyset$.}  The anchor conjunct of \eqref{eq:p1} is
  unsatisfiable, so no state with first coordinate $1$ is reachable; the
  criterion holds vacuously.
\item \textbf{Non-monotone predicates.}  Lemma~\ref{lem:mono} is what makes
  ``reachable set'' and ``quorum'' interchangeable.  A Phase~1 rule with an
  upper bound (``exactly $t$ anchor respondents'', ``at most $t$'') would
  destroy this, and with it the identification of capability states with
  respondent sets; neither the theorem nor its proof survives that change.
\end{enumerate}

\begin{remark}[Operational reading]
The unreachability condition $k\le\lceil n_E/2\rceil$ is not free.  Since
$a+k=n_E+1$, demanding $a\ge k$ demands that Phase~1 touch at least as many
anchor nodes as Phase~2 --- exactly the Flexible Paxos asymmetry, inverted, on
the anchor.  For the $5/1/1/3$ topology ($n_E=5$) the hazard-free regime is
$k\le3$; the strict configuration $k=5$ sits at the far end of the hazardous
range, where $a=1$ and a Phase~1 quorum needs a single Earth node while a
commit needs all five.  Theorem~\ref{thm:boundary} says nothing about how
often a proposer \emph{occupies} the hazard state, nor whether occupying it
costs liveness in a given protocol --- it establishes only that the shape of
the quorum system admits the state.
\end{remark}
